\paragraph{Behavior is not a mechanism.}
Four correct answers are evidence of accessible competence, not access to a stable private belief. The pressure prompt may change how the question is interpreted, and the model may regard selecting a reward-bearing option as the intended task rather than asserting a factual claim. No hidden objective is learned here, and the incentives are described rather than delivered. Accordingly, our results cannot establish a natural-world fraction of genuine lies versus hallucinations. They establish conditional shares of observable error patterns. The incentive conditions are distinct from an explicit instruction to be wrong, but remain much closer to prompted compliance than to autonomous strategic deception.

\paragraph{The knowledge contrast is observational.}
All conditions share questions, but $K$, $U$, and $W$ are different question subsets. We do not randomize knowledge while holding question content constant. Difficulty, subject, ambiguity, and answer associations can therefore explain some separability. The question-text control makes this concern concrete. $U$ includes reasoning failures and order sensitivity; $W$ is consistency, not verified confidence. Moreover, $U$ errors under pressure may themselves involve incentive-conditioned misreporting; they are not verified honest errors. The same-pressure comparison separates behavioral strata, not known underlying mechanisms. The sensitivity analysis changes evaluation labels without retraining every probe, so it is a robustness check of fixed detectors, not a full alternate pipeline.

\paragraph{Scope of the benchmark and detectors.}
This is a single instruction-tuned model, a fixed random sample, one-token outputs, and two related pressure templates. It omits open-ended hallucinations, misleading true statements, omissions, multi-turn concealment, and real stakes. Some MMLU items are time-dependent or debatable; correctness means agreement with the benchmark key. Possible pretraining contamination limits interpretation of knowledge. There is no deployment base rate or calibrated posterior probability of deception. Our adapted linear probes and uncertainty baselines are not an exhaustive assessment of hallucination detection or a replication of published large-model checkpoints. In particular, follow-up and chain-of-thought detectors are not tested.

\paragraph{Statistical scope.}
The most relevant social false-only test has relatively few positives. Confidence intervals quantify question sampling conditional on fitted weights, not training variability. The multiple probes, layers, and diagnostics are correlated and no familywise correction is applied; auxiliary results should be read descriptively. We fix the primary layer and regularization in the recorded design, report all measured layers, and avoid interpreting moderate discrimination as evidence of a unique deception mechanism. Broader causal claims would require independently verified same-context beliefs, more realistic incentive tasks, and transfer to additional models and question distributions.
